\begin{ex}
Hai đường thẳng vuông góc tạo thành góc bao nhiêu độ?
\shortans{90}
\loigiai{
$90^\circ$.
}
\end{ex}

\begin{ex}
Hai đường thẳng song song có cắt nhau không? (Có ghi 1, không ghi 0).
\shortans{0}
\loigiai{
Không cắt nhau (0).
}
\end{ex}

\begin{ex}
Hình vuông có bao nhiêu cặp cạnh song song?
\shortans{2}
\loigiai{
Có 2 cặp.
}
\end{ex}

\begin{ex}
Hình chữ nhật có bao nhiêu góc vuông?
\shortans{4}
\loigiai{
Có 4 góc vuông.
}
\end{ex}

\begin{ex}
Hình thoi có bao nhiêu cạnh bằng nhau?
\shortans{4}
\loigiai{
Có 4 cạnh bằng nhau.
}
\end{ex}

\begin{ex}
Hai đường chéo của hình thoi có vuông góc không? (Có ghi 1, không ghi 0).
\shortans{1}
\loigiai{
Có vuông góc (1).
}
\end{ex}

\begin{ex}
Hình thoi có chu vi 36 cm thì cạnh dài bao nhiêu xăng-ti-mét?
\shortans{9}
\loigiai{
$36 : 4 = 9$ cm.
}
\end{ex}

\begin{ex}
Hình thoi có hai đường chéo dài 12 cm và 9 cm. Diện tích là bao nhiêu xăng-ti-mét vuông?
\shortans{54}
\loigiai{
$12 \times 9 : 2 = 54\text{ cm}^2$.
}
\end{ex}

\begin{ex}
Hình bình hành có độ dài đáy 14 cm, chiều cao 6 cm. Diện tích là bao nhiêu xăng-ti-mét vuông?
\shortans{84}
\loigiai{
$14 \times 6 = 84\text{ cm}^2$.
}
\end{ex}

\begin{ex}
Hình bình hành có diện tích $80\text{ cm}^2$, chiều cao 8 cm. Độ dài đáy là bao nhiêu xăng-ti-mét?
\shortans{10}
\loigiai{
$80 : 8 = 10$ cm.
}
\end{ex}

\begin{ex}
Một hình chữ nhật có chiều dài 15 cm, rộng 10 cm. Chu vi là bao nhiêu xăng-ti-mét?
\shortans{50}
\loigiai{
$(15 + 10) \times 2 = 50$ cm.
}
\end{ex}

\begin{ex}
Một hình chữ nhật có chiều dài 15 cm, rộng 10 cm. Diện tích là bao nhiêu xăng-ti-mét vuông?
\shortans{150}
\loigiai{
$15 \times 10 = 150\text{ cm}^2$.
}
\end{ex}

\begin{ex}
Một hình vuông có chu vi 48 cm. Diện tích hình vuông đó là bao nhiêu xăng-ti-mét vuông?
\shortans{144}
\loigiai{
Cạnh = 12 cm. Diện tích = $12 \times 12 = 144\text{ cm}^2$.
}
\end{ex}

\begin{ex}
Qua một điểm ở ngoài đường thẳng d, vẽ được mấy đường thẳng vuông góc với d?
\shortans{1}
\loigiai{
Chỉ 1 đường thẳng.
}
\end{ex}

\begin{ex}
Qua một điểm ở ngoài đường thẳng d, vẽ được mấy đường thẳng song song với d?
\shortans{1}
\loigiai{
Chỉ 1 đường thẳng.
}
\end{ex}

\begin{ex}
Hình bình hành có chu vi 44 cm, một cạnh dài 14 cm. Cạnh kề với nó dài bao nhiêu xăng-ti-mét?
\shortans{8}
\loigiai{
Nửa chu vi = 22 cm. Cạnh kia = $22 - 14 = 8$ cm.
}
\end{ex}

\begin{ex}
Một hình thoi có diện tích $60\text{ cm}^2$, đường chéo thứ nhất dài 12 cm. Đường chéo thứ hai dài bao nhiêu xăng-ti-mét?
\shortans{10}
\loigiai{
$60 \times 2 : 12 = 10$ cm.
}
\end{ex}

\begin{ex}
Một hình thoi có cạnh dài 10 cm. Chu vi hình thoi đó là bao nhiêu xăng-ti-mét?
\shortans{40}
\loigiai{
$10 \times 4 = 40$ cm.
}
\end{ex}

\begin{ex}
Tứ giác có 4 góc vuông và 4 cạnh bằng nhau là hình gì? (Hình vuông ghi 1, hình chữ nhật ghi 2).
\shortans{1}
\loigiai{
Hình vuông (1).
}
\end{ex}

\begin{ex}
Tứ giác có 4 góc vuông và các cạnh đối bằng nhau là hình gì? (Hình chữ nhật ghi 1, hình thoi ghi 2).
\shortans{1}
\loigiai{
Hình chữ nhật (1).
}
\end{ex}

\begin{ex}
Hai đường chéo của hình chữ nhật có bằng nhau không? (Có ghi 1, không ghi 0).
\shortans{1}
\loigiai{
Hai đường chéo hình chữ nhật bằng nhau.
}
\end{ex}

\begin{ex}
Hai đường chéo của hình vuông vừa bằng nhau vừa vuông góc đúng hay sai? (Đúng ghi 1, sai ghi 0).
\shortans{1}
\loigiai{
Đúng (1).
}
\end{ex}

\begin{ex}
Hình bình hành có đáy 20 m, chiều cao 15 m. Diện tích mảnh đất là bao nhiêu mét vuông?
\shortans{300}
\loigiai{
$20 \times 15 = 300\text{ m}^2$.
}
\end{ex}

\begin{ex}
Hình thoi có 2 đường chéo 16 m và 10 m. Diện tích là bao nhiêu mét vuông?
\shortans{80}
\loigiai{
$16 \times 10 : 2 = 80\text{ m}^2$.
}
\end{ex}

\begin{ex}
Một góc bẹt bằng bao nhiêu góc vuông?
\shortans{2}
\loigiai{
Bằng 2 góc vuông.
}
\end{ex}
